\chapter{Galaxy In The Machine}

Returning to the Computation Center’s CDC 6600 era, the early days in the late sixties and early seventies were special. There was a community of users coexisting together side-by-side in the same environment. Because programming relied entirely on physical punch cards, users had to gather in the Computation Center terminal rooms to punch their code on keypunches. This meant that everyone from aerospace researchers running complex scientific models to undergraduate students debugging basic assignments were working side-by-side. This shared physical space fostered an intellectual commons and  a cross-disciplinary community bound by the limitations of a single, central machine.

Before the proliferation of digital storage and interactive terminals, a program did not exist as a file on a screen. Instead, code was a physical stack of IBM cards. A single line of code, such as a subroutine call passing parameters, was represented by a single card. A complex program required a heavy stack of hundreds of these cards, and it was absolutely critical that they be maintained and fed into the machine in the correct order. To create these cards, users operated a keypunch device, which featured a typewriter-like keyboard used to punch specific patterns of holes into the cards. This process was so central to computing that "Keypunch Operator" was a dedicated job title. The logistics of moving these physical materials also created the role of the "gofer." In high-level research environments, a gofer's primary duty was to physically run stacks of IBM cards and large printouts back and forth between the scientists analyzing the data, the programmers writing the code, the keypunch operators punching the cards, and the computer's input desk. 

At UT, users could submit their decks and retrieve their printouts directly at the Computation Center, or they could walk to a remote job entry site, such as the one housed in Pearce Hall, the Old Law Building, during the early seventies. Perhaps the most extreme example of punch card logistics at UT was the IBM CPC which was donated to the chemistry department by Exxon in 1958. The Card-Programmed Electronic Calculator was an electro-mechanical hybrid that was fundamentally incapable of storing a program's instructions in a digital memory state. For the CPC, the program existed solely as the physical sequence of punched cards. This meant the machine could not execute automated loops on its own. Instead, the human operator had to act as the machine's manual control unit. If an algorithm required an iterative loop to reach a solution, the operator had to stand at the machine and physically re-enter the same card deck over and over again until the calculation was complete.

Once a job was placed in the queue, users often had to wait hours or even days for the computer to process it and return the results on wide-format, green-bar printouts. Because of this builtin delay, a single syntax error on just one card meant restarting the entire logistical loop. This introduced a tremendous amount of cognitive overhead, demanding that programmers plan their code with extreme care before ever approaching the dispatch desk.

While the punch card era was defined by high friction and long waits, it cultivated a unique environment where the university's diverse users were brought together by the shared rituals of batch computing, before the advent of timesharing terminals allowed them to splinter into specialized, decentralized tribes. It was a pioneering era, with a spirit well worth preserving as a historical record and legacy. And there was another force, even more cultural or even instinctual in nature, that tended to bring users together and foster a unique environment. Where there were computers, there were computer games. Games were practically an intrinsic property of these early systems. 

Even the original CDC 6600 checkout engineers famously used the 6600's innovative CRT monitors for early games like Space Wars, Lunar Lander, and Baseball as a way to test the machine. Because it was among the first commercial computers to feature an interactive cathode ray tube display console instead of just glowing lights and teletype text, it became the perfect sandbox for early coders, much like the PDP-1 had a few years earlier with its own single cathode ray tube display. CDC's checkout and maintenance engineers needed a fast, highly visual way to ensure that all parts of the multi-million dollar system, especially the graphics consoles and peripheral processors, were firing correctly under heavy stress. To do this, they programmed a suite of highly advanced, real-time diagnostic games.

While Spacewar was originally coded on the MIT PDP-1 in 1962, the CDC 6600 Space Wars version took full advantage of the supercomputer's relatively immense processing speed. It featured two vector-graphics spaceships maneuvering in real-time, firing torpedoes at each other while being pulled by the gravity of a central star. Lunar Lander was an early, real-time precursor to the text-based and arcade lander games that would explode in popularity in the seventies. Players had to precisely calculate thrust and fuel consumption using the console controls to safely descend a spacecraft onto a jagged vector-graphics moon landscape without crashing. Baseball was a highly unique vector-graphic sports game for the era. A pitcher would throw a pitch, and the batter would have to swing with strict timing to hit the ball out into a digitally rendered diamond. Some historical legal documents from Magnavox patent lawsuits in the seventies point to this exact CDC game as a precursor to early video arcade sports.

CDC engineers openly admitted that the games became the primary incentive for getting the temperamental machines operational. If a newly assembled CDC 6600 could smoothly run Space Wars or Baseball without freezing or crashing, it meant the entire system architecture was completely sound. Because these games utilized the console screens long before commercial video games existed, they can be considered among the first computer games to use graphical monitors.
At UT Austin, this impulse blossomed into a multi-generational family of space simulation games that carried on the Spacewar legacy while melding in the spirit of Star Trek. The story of those early computer games, stretching from the late sixties and the CDC 6600 through the early eighties and the DEC-10, is a dramatic demonstration of how the forces driving computer games could bring users together and keep alive the early pioneering spirit.

\section{National Science Foundation}

The National Science Foundation played a central role in the story of the early days at UT by aggressively funding the Computation Center, including the \$1 million grant that allowed UT to acquire the CDC 6600 in 1966. Initially this computational power was somewhat sequestered  in its dedicated sub-terracian building and oriented towards faculty research projects such as orbital mechanics or structural modeling. The user experience was deep within the batch processing paradigm, with researchers submitting physical stacks of punch cards to a dispatch desk and waiting hours or days for printouts, making computing a high friction and low spontaneity endeavor.

Two landmark initiatives, Project C-BE and Project CONDUIT, were the primary mechanisms through which the NSF sought to break the "Anastasio cycle," a systemic bottleneck where universities possessed powerful mainframes, but a lack of standardized, high-quality educational software prevented their widespread academic adoption. Together, they transformed computing from an intimidating research tool into an accessible, interactive medium for education.

Launched in 1972 and co-directed by chemistry professor Joseph Lagowski and mechanical engineering professor John Allan, Project C-BE focused on local development of computer-augmented curricula. Its overarching philosophy was to use computers to handle routine calculations, grading, and simulations so that human instructors were freed to provide high-level mentoring and direct contact with students. The project ultimately engaged 75 professors and over 4,000 students across 44 curriculum-development projects. It successfully integrated interactive computing into 27 diverse courses, ranging from aerospace engineering and botany to psychology and rhetorical invention in English composition.

Because early teletypes could not print complex math or chemical structures, Project C-BE researchers integrated off-the-shelf analog equipment with the digital system. They wired Kodak Ektagraphic RA-960 random-access 35mm slide projectors to the computer terminals via custom controllers, allowing the mainframe to automatically project high-resolution color graphics on a screen in response to student inputs. To prove the value of these innovations, UT Austin developed the SCRAPE Systematic Control and Evaluation of Academic Programs and Environments model. This framework systematically tracked student achievement, economic efficiency, and program communicability, providing empirical evidence that students using interactive computer-based methods outperformed those in traditional lecture-only environments.
Project CONDUIT was designed to solve the transferability crisis and ensure that software developed at one university wouldn't become obsolete or useless if transferred to another hardware system. Funded by a \$450,000 NSF grant, CONDUIT united five major regional computer networks: UT Austin, Dartmouth College, the University of Iowa, Oregon State University, and the North Carolina Educational Computing Services Center. Through its regional hub, UT Austin helped serve a combined population of 275,000 students across roughly 100 institutions.

CONDUIT became more than a software repository by establishing strict standards. For a program to be distributed nationally, it had to pass both a pedagogical review by disciplinary experts evaluating its educational quality, and a technical review by computer scientists to ensure the code could run on different hardware platforms. The success of this standardization was dramatically proven during a June 1973 bilingual workshop in Mexico City. Latin American educators with no prior programming experience used bilingual guides developed by Project C-BE to successfully execute and interact with standardized programs that had been transported directly from UT Austin and the University of Kansas.

While TAURUS, the Texas Anthropocentric Ubiquitous Responsive User System, was an ingenious software bridge integrated into the locally developed UT-2D operating system, it was ultimately fighting against the CDC 6600's fundamental hardware design. It successfully grafted pseudo-interactive access onto the mainframe, allowing users to connect via remote terminals like the Datapoint 3300 to write and run programs, but the CDC 6600 was engineered by Seymour Cray for raw numeric computation and scientific number crunching, not for the text-heavy, concurrent demands of general-purpose timesharing. During this period, UT researchers encountered Seymour Cray and asked the legendary architect to modify the 6600 hardware so that privileged instructions in user mode would cause an exception rather than a no-op. This was essentially a plea for hardware-level virtualization and improved timesharing support. Cray’s succinct refusal "No, I don’t think so" illustrated the persistent gap between architectural vision and hardware realities. 

The 6600’s architectural limitations actively resisted progress. Text characters were packed ten per 60-bit word using a 6-bit character set, which meant the system natively lacked lower-case letters. Any program trying to manipulate text had to perform complex, inefficient shifting and masking operations. UT Austin even had to adopt a unique 63-character set convention, using null bytes to signal the end of a line. Every running process had to fit entirely within the central memory's strict 17-bit user limit. To keep timesharing somewhat viable, the system had to rely on a two-megaword extended core storage buffer to rapidly swap jobs in and out of central memory. The CDC 6600 offloaded its input and output tasks to ten autonomous PPUs. However, if these PPUs encountered unexpected program states from novice users, they were programmed to hang or spin in an endless loop. This progressively degraded system performance throughout the day until operators performed a scheduled dead-of-night reboot to clear the states. These quirks made it physically and architecturally difficult to support thousands of undergraduate users. 

\section{PDP-10, DEC-10, DEC-20}

A transition away from total reliance on the CDC 6600 and Cyber became a necessity for fully democratizing computing on campus, and the acquisition of the HRC DEC-10 around 1975 was the critical shift. It brought native 7-bit ASCII text handling, virtual memory, and an interrupt-driven architecture designed specifically for high-concurrency timesharing that successfully decoupled the interactive environments of general computing users from rigid, batch-heavy scientific workloads. By transitioning users to remote interactive terminals in the mid seventies, the DEC-10 transformed campus terminal rooms into proto-online social worlds and created a vibrant, collaborative digital culture. This was vividly illustrated by the multi-generational evolution of what began as single-player Star Trek simulations on the CDC 6600 and eventually advanced on the DEC-10 into a real-time, 18-player space simulator that incorporated thrilling interactivity and AI elements. This reflected a deeply connected, multi-user community and a new frontier of digital life at UT Austin. As a contemporary who experienced it noted “I remember seeing Decwar at South Hedland High School, Western Australia in 1981. I wonder how they distributed this from Austin around the world to other PDP-10 sites. We all stood around the few terminals watching and were just in awe, it didn’t seem like a game, but like some science fiction come to life.” 

The social architecture of the Computation Center during this period was defined by an intentional relaxation and outreach, a pedagogical strategy designed to draw young enthusiasts into the STEM pipeline by allowing them direct contact with machines that were otherwise exceptionally rare. During the seventies, the rapid expansion of computing technology was met with significant fear, distrust, and anxiety from the general public, educators, and students. This apprehension stemmed from a combination of unfamiliarity, the perceived threat of dehumanization, and culturally pervasive science fiction. Broad public distrust was captured in a 1971 AFIPS-Time Magazine survey on attitudes toward computing, which revealed strong public anxiety. The study found that approximately one-third of American adults believed the computer was a "thinking machine" capable of thinking for itself. Furthermore, over half of the respondents felt that society was becoming too dependent on computers and that they were changing lives too rapidly, while 15\% believed computers actually made their lives worse.

These fears were heavily influenced by popular media and science fiction stories of mechanical servants and electronic brains, epitomized by HAL, the murderous computer from the 1968 film 2001: A Space Odyssey. This widespread computer illiteracy was viewed as a societal threat. As Dartmouth College President John Kemeny warned in a 1976 CBS interview, graduating "computer illiterates" was highly dangerous because it meant "a small number of computer experts will make decisions for us as to how computers will be used".

Within schools and universities, teachers and administrators frequently viewed computers with deep suspicion. Educators tended to "dislike and distrust computers because they do not understand them". Because the technology was highly unfamiliar, teachers often saw computers as a direct threat. At institutions like St. Stephen's Episcopal School in Austin, early computer integration was met with student enthusiasm but very little faculty involvement because teachers harbored "a fear of being replaced by the machine". Teachers worried that technology would threaten their status with children and cause them to lose control over what was taught and how pupil activities were scheduled. There was a prevalent myth that educational technology would inherently bring about "regimentation and a loss of individualization". Some feared that interacting with a terminal would be a depersonalizing experience that stripped away human contact.

For students, the fear of computers was often tied to math anxiety. Early mainframe computers were seen purely as complex "number crunchers," which induced a "fear of handling numbers" and general frustration, particularly among students in the behavioral sciences and humanities. To combat these deeply ingrained fears, universities dedicated significant resources to changing the culture around computing in the seventies. At UT Austin, the NSF-funded Project C-BE was specifically designed to shift the student experience away from "computational intimidation" and toward "interactive empowerment". By using simplified programming languages like OMNITAB, educators sought to remove the "frustration, fear, and boredom" of complex calculations. To overcome "faculty illiteracy" and resistance, universities ran intensive training workshops designed to instill confidence and prove to teachers that they were in control of the machines, not the other way around. In 1976, Dartmouth College hosted a national conference for over 200 administrators specifically aimed at "demystifying the computer".

During the late sixties and seventies, the general public and countercultural activists largely viewed mainframe computers as monolithic threats, instruments of centralized power, military surveillance, and corporate bureaucracy closely tied to the Cold War. The DEC-10 actively subverted this image, transforming "Big Iron" into a shared canvas for personal and collective liberation. The dominant computing paradigm had been batch processing, where users handed physical punch cards to a specialized "priesthood" of machine operators who controlled all access to the hardware. The DEC-10 disrupted this centralized hierarchy by popularizing interactive timesharing. 

DEC’s 36-bit mainframe systems reflected an evolution stage in computer engineering, system software design, and industrial marketing. It was an ecosystem founded on the Programmed Data Processor-10 or PDP-10, an overarching architecture introduced in 1966 as the direct descendant of the PDP-6, which had been DEC’s first effort to compete with IBM’s 7090 series 36-bit machines in the large computer market. Defined by its 36-bit word length and instruction set, the PDP-10 was not a single physical machine but a foundational computing architecture implemented across successive hardware processor generations, ranging from the early discrete-transistor KA10 and KI10 to the later, microprogrammed KL10 and KS10 processors. It served as the primary hardware substrate for the pioneering timesharing experiments and early artificial intelligence research at institutions such as MIT and Stanford.

As DEC sought to expand its market presence beyond academic and scientific research laboratories into the commercial business sector, the company instituted a significant rebranding strategy. In 1971, systems built upon the PDP-10 architecture were commercialized under the name DECsystem-10, commonly abbreviated as DEC-10. Mechanically, a DEC-10 was a hardware assembly utilizing KA10 or KI10 processors bundled with DEC’s proprietary Tops-10 operating system. This configuration optimized raw hardware efficiency and supported robust, multi-user timesharing environments, cementing the platform's reputation as a dependable workhorse for university computing centers and commercial service bureaus throughout the seventies.

The introduction of the DECSYSTEM-20 or DEC-20 in 1976 represented both a hardware modernization and a fundamental shift in user-system interaction. Utilizing the advanced KL10 and KS10 processors, the DEC-20 architecture was distinguished primarily by its integration of the TOPS-20 operating system. TOPS-20 evolved from TENEX, a third-party operating system developed by Bolt, Beranek and Newman that featured sophisticated paged virtual memory architecture. While the underlying 36-bit instruction set remained compatible with its predecessors, the DEC-20 provided an entirely modernized user experience characterized by an advanced command-line interface featuring interactive command completion and integrated help facilities, marking a significant milestone in the evolution of user-centric operating systems.

By giving users direct, real-time access to the processor via remote terminals, the PDP-10 eliminated the bureaucratic middlemen. This shift strongly aligned with the countercultural emphasis on personal autonomy, direct experience, and the democratization of tools. Activists like Ted Nelson explicitly championed timesharing systems in his 1974 manifesto Computer Lib, arguing against the restrictive "priesthood" and urging readers to chant, "Computer power to the people!". The countercultural ethos of the PDP-10 was deeply embedded in its operating systems, which were often designed by the hackers themselves. 

At the MIT Artificial Intelligence Laboratory, hackers built the Incompatible Timesharing System as a protest against structured, secure corporate operating systems. Reflecting the belief that information should be free, ITS was built with no passwords, no file protection, and no central security controls. Logging in was merely a courtesy, and anyone on the Arpanet could log in as an anonymous tourist, creating an environment of radical openness and peer-to-peer collaboration. At Stanford’s Artificial Intelligence Laboratory, the PDP-10 ran the WAITS operating system in a semicircular wooden building nestled in the Santa Cruz Mountains. The SAIL community blended academic research with a bohemian lifestyle, featuring a clothing-optional sauna, people riding horses to work, and rooms named after locations in J.R.R. Tolkien’s The Lord of the Rings.

The PDP-10 proved to be an engine for digital play, which served as a powerful tool to demystify computing. Stewart Brand, the countercultural publisher of the Whole Earth Catalog, recognized this when he visited Stanford's SAIL to observe programmers playing Spacewar on their PDP-10 screens. In his highly influential 1972 Rolling Stone article, Brand framed these hackers as a new creative elite akin to Ken Kesey's Merry Pranksters. He famously declared that "the computer itself is a new LSD," framing interactive computing as a consciousness-expanding tool that would undermine bureaucracies and bring about a more flexible, playful social world.

Perhaps the greatest paradox of this era was that the anti-authoritarian culture of the PDP-10 labs was almost entirely funded by the United States military ARPA/DARPA research organizations to build robust command-and-control infrastructure. The hackers engaged in a form of technological reclamation. They took the primary instruments of Cold War defense research and repurposed them to build a decentralized, collaborative digital playground. Instead of just weapon systems, they used the military's PDP-10s to invent network email, text-adventure games, and chat environments.

The PDP-10 was a unique technological platform where hardware architecture and software design became direct expressions of countercultural values. By replacing the rigid corporate batch-processing model with an interactive, unmediated environment, it established the social and ethical frameworks that would later produce the free software movement and the collaborative structure of the modern internet. It reflected new horizons and new frontiers opening up at UT Austin in the seventies, and the beginnings of a unique story that is still alive today, fifty years later.
